% Page 042

\seite{42}

\abschnitt{Revision}
\pstart
Am 15.\ Januar wurde die Schule durch\ob
den Herrn Kreisschulinspektor von Bitburg\ob
revidiert.

\pend
\abschnitt{Weihnachtsfeier}
\pstart

Zum Geburtstag Seiner Majestät des Kaisers\ob
wurde eine gemeinsame Feier mit den Schulen\ob
von Seffern und Schleid abgehalten. Die\ob
Kinder trugen Lieder vor, und Herr Pfarrer\ob
von Schleid hielt eine Ansprache. Kinder\ob
aus den umliegenden Ortschaften sangen\ob
gleichfalls.

\pend
\abschnitt{Revision}
\pstart

Am 25.\ Februar wurde die Schule durch\ob
den Herrn Kreisschulinspektor Gallin von\ob
Bitburg revidiert.

\pend
\abschnitt{Osterprüfung}
\pstart

Am 18.\ März fand durch den Herrn\ob
Ortsschulinspektor Pfarrer Schmitt die\ob
übliche Osterprüfung statt.

\pend
\jahresueberschrift{Das Schuljahr 1906-07}
\pstart

Es wurden entlassen einige Knaben und\ob
1 Mädchen. Die Knaben widmen sich der\ob
Landwirthschaft. Aufgenommen wurden\ob
4 Knaben und 3 Mädchen, sodaß die Schule\ob
nunmehr 51 Kinder zählt. Alle Kinder sind\ob
katholisch und willfährig; ein Knabe ist\ob
nicht ganz willfährig.

Am 15.\ Mai ward der im Jahre 1839 ge\ohb
borene Lehrer hier\unsicher, der an hiesiger\ob
Stelle 15 Jahre wirkte, verstorben.

\pend
\abschnitt{den 27./3. 1906}
\pstart

Schmitt, Pfarrer. {[Unterschrift]}

\pend
\abschnitt{Vertretung}
\pstart

Am 1.\ August wurde eine neue Vertretungs\ohb
regelung getroffen, wofür 60 M gezahlt\ob
wurden.

\pend
\abschnitt{Ausfall des Unterrichts}
\pstart

Wegen großer Hitze im August und Sep\ohb
tember fiel der Nachmittagsunterricht\ob
aus.
\pend
